\documentclass[a4paper,11pt]{article}
\def\email#1{{\tt#1}}

\def\N{\mbox{\rm{I}\hspace{-.15em}\rm{N}}}
\def\Z{\mbox{\rm{Z}\hspace{-.28em}\rm{Z}}}
\def\RR{\mbox{\rm{I}\hspace{-.15em}\rm{R}}}
\def\C{\mbox{\rm{l}\hspace{-.50em}\rm{C}}}
\def\Q{\mbox{\rm{l}\hspace{-.50em}\rm{Q}}}
\def\F{\mbox{\rm{I}\hspace{-.20em}\rm{F}}}

\usepackage{amssymb,amsmath}
\usepackage{url}
\usepackage{fullpage}
\begin{document}

\newtheorem{theo}{Theorem}
\newtheorem{coro}{Corollary}[theo]
\newtheorem{lemme}{Lemma}
\renewcommand{\thecoro}{\thetheo.\arabic{coro}}



\parindent=0cm


\section*{\center IN100 Initiation au Python\\
Contrôle continu 2}

\vspace{1 cm}

\begin{table}[h!]
\centering
\renewcommand{\arraystretch}{1.4} % augmente la hauteur des lignes
\begin{tabular}{|p{10cm}|p{3cm}|p{2cm}|}
\hline
\textbf{Étudiant :} & \textbf{Heure Début }  & \textbf{Note (/5) }  \\
& & \\
\hline
\textbf{Correcteur :} & \textbf{Heure Fin } & \\
& & \\
\hline

\hline
\end{tabular}
\end{table}

\vspace{-0.2 cm}
\subsection*{Consignes:} Vous avez 20 minutes pour coder la solution. Seules les types et les instructions vues en cours sont autorisées. L'usage d'une fonction 
venant d'une bibliothèque qui n'a pas été vue en cours ne sera pris en compte. 

\vspace{-0.2 cm}
\subsection*{Sujet}

Une pile est une structure de données basée sur le principe que le dernier élément ajouté dans la pile est le premier à en être retiré.
Nous l'utiliserons pour gérer une pile de livres à ranger sur une étagère. Chaque livre est représenté par un entier correspondant à son nombre de pages.

\begin{itemize}
\item[1.] \'Ecrire une fonction \texttt{ajoute\_livre} qui prend en argument une liste et un entier \texttt{nb\_pages}, et qui ajoute le livre dans la liste (en haut de la pile, donc en fin de liste).
Testez votre fonction en ajoutant les livres suivants dans la pile (dans cet ordre) :
\vspace{-0.2 cm}
\begin{verbatim}
150
200
120
80
220
\end{verbatim}
\vspace{-0.2 cm}
\item[2.] \'Ecrire une fonction \texttt{ranger\_livre} prend en argument la liste qui donne la pile, enlève le dernier élément ajouté (le livre en haut de la pile),
et renvoie le nombre de pages de ce livre.
\item[3.] \'Ecrire une fonction \texttt{ranger\_pile\_livres} qui prend en argument  la pile de livres, traite tous les livres dans la pile jusqu’à ce qu’elle soit vide,
et renvoie le total de pages de tous les livres rangés. 
\end{itemize}
\vspace{-0.2 cm}
\begin{table}[h!]
\centering
\renewcommand{\arraystretch}{1.4} % augmente la hauteur des lignes
\begin{tabular}{|p{15cm}|}
\hline
\textbf{Remarques étudiant :}    \\
 \\
\hline
\textbf{Remarques correcteur :}  \\
 \\
\hline
\end{tabular}
\end{table}

\newpage

\vspace{-0.2 cm}

\section*{\center IN100 Initiation au Python\\
Contrôle continu 2}

\vspace{-0.3 cm}

\begin{table}[h!]
\centering
\renewcommand{\arraystretch}{1.4} % augmente la hauteur des lignes
\begin{tabular}{|p{10cm}|p{3cm}|p{2cm}|}
\hline
\textbf{Étudiant :} & \textbf{Heure Début }  & \textbf{Note (/5) }  \\
& & \\
\hline
\textbf{Correcteur :} & \textbf{Heure Fin } & \\
& & \\
\hline

\hline
\end{tabular}
\end{table}
\vspace{-0.4 cm}
\paragraph*{Consignes:} Vous avez 20 minutes pour coder la solution. Seules les types et les instructions vues en cours sont autorisées. Tout usage d'une fonction 
venant d'une bibliothèque qui n'a pas été vue en cours ne sera pris en compte. 

\vspace{-0.2 cm}
\paragraph*{Sujet:}

Un \textbf{code produit} est composé de 12 chiffres : 
\vspace{-0.2 cm}
\begin{itemize}
    \item Les 11 premiers chiffres correspondent aux informations sur le produit.
    \item Le dernier chiffre représente la \textbf{clé de contrôle}, calculée à partir des 11 premiers chiffres.
\end{itemize}
\vspace{-0.2 cm}
La clé de contrôle est calculée ainsi :
\vspace{-0.2 cm}

\[
\text{clé} = (\text{somme des 11 premiers chiffres}) \bmod 10
\]
\vspace{-0.2 cm}

\textbf{Exemple :}  
Si les 11 premiers chiffres sont \texttt{12345678901}, alors :  
\vspace{-0.2 cm}

\[
\text{clé} = (1+2+3+4+5+6+7+8+9+0+1) \mod 10 = 46 \mod 10 = 6
\]
\vspace{-0.3 cm}

Le code complet sera donc \texttt{123456789016}, où 6 est la clé calculée.

\vspace{-0.3 cm}


\begin{enumerate}
  \item[1.] Écrire une fonction qui demande à l'utilisateur de saisir un \textbf{code produit} sous forme de chaîne de caractères.
 \item[2.] Écrire une fonction Python  \texttt{verifier\_code\_produit}  qui vérifie qu'un code produit est valide. 
 \item[3.] \'Ecrire une fonction qui génère un code produit,  les premières 11 caractères étant générés aléatoirement. 
 \end{enumerate}
 \vspace{-0.2 cm}

Pour générer aléatoirement des nombres entiers, nous rappelons qu'on peut utiliser la fonction \texttt{random.randint(a,b)} qui renvoie un nombre entier entre $a$ et $b$. 

\vspace{0.2 cm}
Exemple d'usage:
\vspace{-0.2 cm}
\begin{verbatim}
import random

print(random.randint(2,15))
\end{verbatim}

\vspace{-0.4 cm}
\begin{table}[h!]
\centering
\renewcommand{\arraystretch}{1.4} % augmente la hauteur des lignes
\begin{tabular}{|p{15cm}|}
\hline
\textbf{Remarques étudiant :}    \\
 \\
\hline
\textbf{Remarques correcteur :}  \\
 \\
\hline
\end{tabular}
\end{table}


\newpage


\section*{\center IN100 Initiation au Python\\
Contrôle continu 2}


\begin{table}[h!]
\centering
\renewcommand{\arraystretch}{1.4} % augmente la hauteur des lignes
\begin{tabular}{|p{10cm}|p{3cm}|p{2cm}|}
\hline
\textbf{Étudiant :} & \textbf{Heure Début }  & \textbf{Note (/5) }  \\
& & \\
\hline
\textbf{Correcteur :} & \textbf{Heure Fin } & \\
& & \\
\hline

\hline
\end{tabular}
\end{table}

\paragraph*{Consignes:} Vous avez 20 minutes pour coder la solution. Seules les types et les instructions vues en cours sont autorisées. Tout usage d'une fonction 
venant d'une bibliothèque qui n'a pas été vue en cours ne sera pris en compte. 

\paragraph*{Sujet:} On souhaite créer un petit système de chiffrement simple basé sur une permutation de l’alphabet.

\begin{itemize}
\item[1.] Écrire une fonction \texttt{permutation\_alphabet} qui 
crée une permutation aléatoire de l’alphabet. Pour cela vous utiliserez la fonction
\texttt{shuffle} du module \texttt{random}. 
En partant d'une liste \texttt{alphabet} contenant les caractères \texttt{'a'} à l'indice 0, \texttt{'b'} à l'indice 1, \texttt{'c'} à l'indice 2 , \texttt{'d'} à l'indice 4  \ldots 
on peut faire: 
\begin{verbatim}
import random

random.shuffle(alphabet)
\end{verbatim}
\item[2.]  Écrire une fonction \texttt{chiffrer\_mot} 
qui prend en argument le mot à chiffrer  et la liste donnant la permutation et 
qui renvoie le mot chiffré, en remplaçant chaque lettre de l'alphabet par la lettre se trouvant au même indice dans la liste permutée. \\
\vspace{-0.1 cm}

Exemple: 
\vspace{-0.2 cm}
\begin{verbatim}
perm=['u', 'n', 'o', 'd', 'k', 'c', 'g', 'm', 'f', 'h', 'x', 't', 'w', 'j',
 'e', 's', 'v', 'y', 'l', 'r', 'z', 'a', 'p', 'b', 'q', 'i']
>>> chiffrer_mot("bonjour", perm)
'nej...'
\end{verbatim}

\item[3.]  Écrire une fonction \texttt{dechiffrer\_mot} 
qui prend en argument le mot chiffré et la liste donnant la permutation et 
qui renvoie le mot en clair. 
\end{itemize}


\vspace{-0.4 cm}

\begin{table}[h!]
\centering
\renewcommand{\arraystretch}{1.4} % augmente la hauteur des lignes
\begin{tabular}{|p{15cm}|}
\hline
\textbf{Remarques étudiant :}    \\
 \\
\hline
\textbf{Remarques correcteur :}  \\
 \\
\hline
\end{tabular}
\end{table}


\newpage

\section*{\center IN100 Initiation au Python\\
Contrôle continu 2}


\begin{table}[h!]
\centering
\renewcommand{\arraystretch}{1.4} % augmente la hauteur des lignes
\begin{tabular}{|p{10cm}|p{3cm}|p{2cm}|}
\hline
\textbf{Étudiant :} & \textbf{Heure Début }  & \textbf{Note (/5) }  \\
& & \\
\hline
\textbf{Correcteur :} & \textbf{Heure Fin } & \\
& & \\
\hline

\hline
\end{tabular}
\end{table}

\vspace*{-0.4 cm}
\subsection*{Consignes:} Vous avez 20 minutes pour coder la solution. Seules les types et les instructions vues en cours sont autorisées. Tout usage d'une fonction 
venant d'une bibliothèque qui n'a pas été vue en cours ne sera pris en compte. 

\vspace*{-0.2 cm}

\subsection*{Sujet:}   


La suite de Fibonacci est définie par la relation suivante :
\[
F_0 = 0, \quad F_1 = 1, \quad \text{et pour tout } n \ge 2, \quad F_n = F_{n-1} + F_{n-2}.
\]

Autrement dit, chaque terme est la somme des deux précédents.

\begin{itemize}
\item[1.] Écrire une fonction \texttt{fibonacci} qui prend en entrée un entier $n$ et qui calcule le $n$-ième nombre de Fibonnaci. 
\item[2.] Écrire une fonction  \texttt{est\_fibonacci}  qui prend en entrée un entier $x$. La fonction devra  générer la suite jusqu’à dépasser la valeur \texttt{x}, puis vérifier si \texttt{x} y apparaît.
\end{itemize}



\begin{table}[h!]
\centering
\renewcommand{\arraystretch}{1.4} % augmente la hauteur des lignes
\begin{tabular}{|p{15cm}|}
\hline
\textbf{Remarques étudiant :}    \\
 \\
 \\
\hline
\textbf{Remarques correcteur :}  \\
 \\
 \\
\hline
\end{tabular}
\end{table}

\newpage

\section*{\center IN100 Initiation au Python\\
Contrôle continu 2}


\begin{table}[h!]
\centering
\renewcommand{\arraystretch}{1.4} % augmente la hauteur des lignes
\begin{tabular}{|p{10cm}|p{3cm}|p{2cm}|}
\hline
\textbf{Étudiant :} & \textbf{Heure Début }  & \textbf{Note (/5) }  \\
& & \\
\hline
\textbf{Correcteur :} & \textbf{Heure Fin } & \\
& & \\
\hline

\hline
\end{tabular}
\end{table}
\vspace{-0.6 cm}
\subsection*{Consignes:} Vous avez 20 minutes pour coder la solution. Seules les types et les instructions vues en cours sont autorisées. Tout usage d'une fonction 
venant d'une bibliothèque qui n'a pas été vue en cours ne sera pris en compte. 

\subsection*{Sujet:}

Dans une salle de cinéma, les places sont représentées par une \textbf{matrice} de taille $n \times n$.  

\begin{itemize}
\item Chaque ligne représente une \textbf{rangée} de sièges.  
\item Chaque colonne représente une \textbf{place} dans la rangée.  
\end{itemize}
Une valeur \texttt{1} indique qu'une place est \textbf{occupée}, \texttt{0} qu'elle est \textbf{libre}.
On écrit un programme qui permettra au gestionnaire au guichet de vendre les tickets. 

\begin{itemize}
\item[1.] \'Ecrire une fonction prend en entrée la matrice d'une salle de théâtre, et un entier représentant le nombre de tickets que l'utilisateur souhaite 
acheter et renvoie une liste avec les places qui ont été attribuées. La liste contiendra des tuples d'entiers \texttt{(i,j)} donnant le numéro de la rangée et le numéro de la place sur la rangée. 
\item[2.] \'Ecrire une fonction qui, à la fin des ventes, comptera le nombre de places vendues. Cette fonction prendra en entrée la matrice de la salle. 
\end{itemize}


%\vspace*{-0.3 cm}
\begin{table}[h!]
\centering
\renewcommand{\arraystretch}{1.4} % augmente la hauteur des lignes
\begin{tabular}{|p{15cm}|}
\hline
\textbf{Remarques étudiant :}    \\
 \\
 \\
\hline
\textbf{Remarques correcteur :}  \\
 \\
 \\
\hline
\end{tabular}
\end{table}





\end{document}
